\BLOCK{include '_preamble.tex'}

\begin{document}
\maketitle

%% ---- POCETAK PRORACUNA ----

\subsection{Ulazni parametri}

\begin{table}[H]
\centering
\begin{tabular}{|c|l|c|c|}
\hline
Oznaka & Opis & Vrednost & Jedinica \\
\hline
$N_{Ed}$ & Normalna sila pritiska & \VAR{m.NEd|n} & kN \\
\hline
$b$ & Debljina zida & \VAR{(m.b*100)|n(1)} & cm \\
\hline
$h$ & Dužina preseka zida & \VAR{(m.h*100)|n(1)} & cm \\
\hline
$f_{ck}$ & Karakteristična čvrstoća betona & \VAR{m.fck|n(0)} & MPa \\
\hline
$q_0$ & Osnovna vrednost faktora ponašanja & \VAR{m.q0|n} & - \\
\hline
$h_0$ & Dužina utegnutog elementa & \VAR{(m.h_0*100)|n(1)} & cm \\
\hline
$O$ & Obim uzengija za pridržavanje & \VAR{(m.obim_uzengija*100)|n(1)} & cm \\
\hline
$s$ & Vertikalni razmak uzengija & \VAR{(m.s*100)|n(1)} & cm \\
\hline
- & Tip armature & \VAR{m.tip_armature} & - \\
\hline
$\phi_w$ & Prečnik uzengije & \VAR{precnik} & mm \\
\hline
$T$ & Period oscilovanja & \VAR{m.T|n} & s \\
\hline
$M_{Ed}$ & Seizmički moment & \VAR{m.MEd|n} & kNm \\
\hline
$M_{Rd}$ & Moment nosivosti & \VAR{m.MRd|n} & kNm \\
\hline
$H_s$ & Čista spratna visina & \VAR{m.Hs|n} & m \\
\hline
\BLOCK{set pola = (bi|length + 1) // 2}
\BLOCK{if bi|length > 1}
%# Razmaci b_i u dva reda tabele (bez linije između), da stanu u širinu stranice
$b_i$ & Razmaci pridržanih šipki & \VAR{bi[:pola]|map('n', 1)|join(', ')}, & cm \\
 & & \VAR{bi[pola:]|map('n', 1)|join(', ')} & \\
\BLOCK{else}
$b_i$ & Razmaci pridržanih šipki & \VAR{bi|map('n', 1)|join(', ')} & cm \\
\BLOCK{endif}
\hline
\end{tabular}
\caption{Ulazni parametri zida \VAR{naziv|tex}}
\end{table}

\begin{figure}[H]
\centering
\includegraphics[height=0.42\textheight]{\VAR{slike.presek}}
\caption{Presek zida sa utegnutim krajevima}
\end{figure}

\subsection{Materijal i geometrija}

\begin{align*}
f_{cd} &= \frac{0.85 \cdot f_{ck}}{1.5} = \frac{0.85 \cdot \VAR{m.fck|n(0)}}{1.5} = \VAR{(m.fcd/1000)|n} \ \mathrm{MPa} \\
f_{yd} &= \frac{f_{yk}}{\gamma_s} = \frac{500}{1.15} = \VAR{(m.fyd/1000)|n} \ \mathrm{MPa} \\
b_0 &= b - 5 = \VAR{(m.b*100)|n(1)} - 5 = \VAR{(m.b_0*100)|n(1)} \ \mathrm{cm} \\
\omega_v &= \rho_{v,min} \cdot \frac{f_{yd}}{f_{cd}} = 0.00201 \cdot \frac{\VAR{(m.fyd/1000)|n}}{\VAR{(m.fcd/1000)|n}} = \VAR{m.omega_v|n(4)} \\
A_{sw} &= \frac{\phi_w^2 \cdot \pi}{4} = \VAR{(m.a1_uz*10000)|n(3)} \ \mathrm{cm^2}
\end{align*}

\subsection{Normalizovana aksijalna sila}

\begin{align*}
\nu_{d} &= \frac{N_{Ed}}{b \cdot h \cdot f_{cd}} = \frac{\VAR{m.NEd|n}}{\VAR{m.b|n(3)} \cdot \VAR{m.h|n(3)} \cdot \VAR{(m.fcd/1000)|n} \cdot 10^3} = \VAR{m.ni|n(3)}
\end{align*}

Uslov potrebe za utezanjem krajeva zida:

\begin{align*}
\nu_d &> 0.15 \\
\VAR{m.ni|n(3)} &> 0.15 \quad \VAR{uslov(m.ni > 0.15)}
\end{align*}

\BLOCK{if d.slucaj == 'nije_potrebno'}
\textbf{Zaključak:} nije potrebno utezanje krajeva zida \VAR{naziv|tex}.
\BLOCK{else}
Kontrola maksimalne normalizovane sile za seizmički zid:

\begin{align*}
\nu_d &\leq 0.40 \\
\VAR{m.ni|n(3)} &\leq 0.40 \quad \VAR{kontrola(m.ni <= 0.4)}
\end{align*}

\BLOCK{if d.slucaj == 'prekoracena_sila'}
\textbf{Zaključak:} prekoračena je maksimalna normalizovana sila za seizmički zid \VAR{naziv|tex}; dalji proračun utezanja se ne sprovodi.
\BLOCK{else}

\subsection{Faktor duktilnosti krivine $\mu_\phi$}

Prema EN 1998-1, 5.4.3.4.2(2), u izrazima (5.4) i (5.5) iz 5.2.3.4(3) osnovna vrednost $q_0$ se zamenjuje proizvodom $q_0$ i odnosa $M_{Ed}/M_{Rd}$ u dnu zida:

\begin{align*}
q_0 \cdot \frac{M_{Ed}}{M_{Rd}} &= \VAR{m.q0|n} \cdot \frac{\VAR{m.MEd|n}}{\VAR{m.MRd|n}} = \VAR{(m.q0*m.MEd/m.MRd)|n(3)}
\end{align*}

\BLOCK{if d.predimenzionisan}
Kako je ovaj proizvod manji od 1, usvaja se $q_0' = 1$.
\BLOCK{else}
Usvaja se $q_0' = \VAR{d.q0p|n(3)}$.
\BLOCK{endif}

Za armaturu \VAR{m.tip_armature} koeficijent prema 5.2.3.4(4) iznosi $k = \VAR{k_arm|n(1)}$.

\BLOCK{if m.T >= Tc}
Kako je $T_1 = \VAR{m.T|n} \ \mathrm{s} \geq T_C = \VAR{Tc|n} \ \mathrm{s}$, izraz (5.4):

\begin{align*}
\mu_\phi &= k \cdot (2 \cdot q_0' - 1) = \VAR{k_arm|n(1)} \cdot (2 \cdot \VAR{d.q0p|n(3)} - 1) = \VAR{d.mfi|n(3)}
\end{align*}
\BLOCK{else}
Kako je $T_1 = \VAR{m.T|n} \ \mathrm{s} < T_C = \VAR{Tc|n} \ \mathrm{s}$, izraz (5.5):

\begin{align*}
\mu_\phi &= k \cdot \left( 1 + 2 \cdot (q_0' - 1) \cdot \frac{T_C}{T_1} \right) = \VAR{k_arm|n(1)} \cdot \left( 1 + 2 \cdot (\VAR{d.q0p|n(3)} - 1) \cdot \frac{\VAR{Tc|n}}{\VAR{m.T|n}} \right) = \VAR{d.mfi|n(3)}
\end{align*}
\BLOCK{endif}

\subsection{Visina pritisnute zone i potreban procenat utezanja}

\begin{align*}
x_u &= (\nu_d + \omega_v) \cdot h \cdot \frac{b}{b_0} = (\VAR{m.ni|n(3)} + \VAR{m.omega_v|n(4)}) \cdot \VAR{(m.h*100)|n(1)} \cdot \frac{\VAR{(m.b*100)|n(1)}}{\VAR{(m.b_0*100)|n(1)}} = \VAR{(d.xu*100)|n(1)} \ \mathrm{cm} \\
\alpha \cdot \omega_{wd,req} &= 30 \cdot \mu_\phi \cdot (\nu_d + \omega_v) \cdot \varepsilon_{sy,d} \cdot \frac{b}{b_0} - 0.035 \\
&= 30 \cdot \VAR{d.mfi|n(3)} \cdot (\VAR{m.ni|n(3)} + \VAR{m.omega_v|n(4)}) \cdot \VAR{esyd|n(6)} \cdot \frac{\VAR{(m.b*100)|n(1)}}{\VAR{(m.b_0*100)|n(1)}} - 0.035 = \VAR{d.alfa_omegad_req|n(4)}
\end{align*}

\subsection{Minimalna dužina utezanja}

\begin{align*}
\varepsilon_{cu2,c} &= 0.0035 + 0.1 \cdot \alpha \cdot \omega_{wd,req} = 0.0035 + 0.1 \cdot \VAR{d.alfa_omegad_req|n(4)} = \VAR{d.epsilon_cu2_c|n(5)} \\
l_{c,min} &= x_u \cdot \left(1 - \frac{\varepsilon_{cu2}}{\varepsilon_{cu2,c}}\right) = \VAR{(d.xu*100)|n(1)} \cdot \left(1 - \frac{0.0035}{\VAR{d.epsilon_cu2_c|n(5)}}\right) = \VAR{(d.lc_mreq_racunsko*100)|n(1)} \ \mathrm{cm}
\end{align*}

Uslov minimalne dužine utezanja:

\begin{align*}
l_{c,min} &\geq \max(0.15 \cdot h;\ 1.5 \cdot b) = \max(\VAR{(0.15*m.h*100)|n(1)};\ \VAR{(1.5*m.b*100)|n(1)}) = \VAR{([0.15*m.h, 1.5*m.b]|max*100)|n(1)} \ \mathrm{cm} \\
\VAR{(d.lc_mreq_racunsko*100)|n(1)} &\geq \VAR{([0.15*m.h, 1.5*m.b]|max*100)|n(1)} \quad \VAR{uslov(d.lc_mreq_racunsko >= [0.15*m.h, 1.5*m.b]|max)}
\end{align*}

\BLOCK{if d.lc_mreq_racunsko < [0.15*m.h, 1.5*m.b]|max}
Usvaja se minimalna dužina utezanja $l_{c,min} = \VAR{(d.lc_mreq*100)|n(1)} \ \mathrm{cm}$.
\BLOCK{else}
Minimalna dužina utezanja je $l_{c,min} = \VAR{(d.lc_mreq*100)|n(1)} \ \mathrm{cm}$.
\BLOCK{endif}

\subsection{Koeficijenti efikasnosti utezanja}

\begin{align*}
\alpha_n &= 1 - \frac{\sum b_i^2}{6 \cdot b_0 \cdot h_0} = 1 - \frac{\VAR{bi|map('n', 1)|join('^2 + ')}^2}{6 \cdot \VAR{(m.b_0*100)|n(1)} \cdot \VAR{(m.h_0*100)|n(1)}} = \VAR{d.alfa_n|n(4)} \\
\alpha_s &= \left(1 - \frac{s}{2 \cdot b_0}\right) \cdot \left(1 - \frac{s}{2 \cdot h_0}\right) = \left(1 - \frac{\VAR{(m.s*100)|n(1)}}{2 \cdot \VAR{(m.b_0*100)|n(1)}}\right) \cdot \left(1 - \frac{\VAR{(m.s*100)|n(1)}}{2 \cdot \VAR{(m.h_0*100)|n(1)}}\right) = \VAR{d.alfa_s|n(4)} \\
\alpha &= \alpha_n \cdot \alpha_s = \VAR{d.alfa_n|n(4)} \cdot \VAR{d.alfa_s|n(4)} = \VAR{d.alfa|n(4)}
\end{align*}

\subsection{Obezbeđeni mehanički zapreminski procenat utezanja}

\begin{align*}
V_{sw} &= A_{sw} \cdot O = \VAR{(m.a1_uz*10000)|n(3)} \cdot \VAR{(m.obim_uzengija*100)|n(1)} = \VAR{(d.Vsw*1e6)|n(2)} \ \mathrm{cm^3} \\
V_{co} &= b_0 \cdot h_0 \cdot s = \VAR{(m.b_0*100)|n(1)} \cdot \VAR{(m.h_0*100)|n(1)} \cdot \VAR{(m.s*100)|n(1)} = \VAR{(d.Vco*1e6)|n(1)} \ \mathrm{cm^3} \\
\omega_{wd,prov} &= \frac{V_{sw} \cdot f_{yd}}{V_{co} \cdot f_{cd}} = \frac{\VAR{(d.Vsw*1e6)|n(2)} \cdot \VAR{(m.fyd/1000)|n}}{\VAR{(d.Vco*1e6)|n(1)} \cdot \VAR{(m.fcd/1000)|n}} = \VAR{d.omega_d_prov|n(4)} \\
\omega_{wd,req} &= \max\left(0.08;\ \frac{\alpha \cdot \omega_{wd,req}}{\alpha}\right) = \max\left(0.08;\ \frac{\VAR{d.alfa_omegad_req|n(4)}}{\VAR{d.alfa|n(4)}}\right) = \VAR{d.limit_omega|n(4)}
\end{align*}

Kontrola procenta utezanja:

\begin{align*}
\omega_{wd,prov} &\geq \omega_{wd,req} \\
\VAR{d.omega_d_prov|n(4)} &\geq \VAR{d.limit_omega|n(4)} \quad \VAR{kontrola(d.dovoljno_omega)}
\end{align*}

\subsection{Potrebna dužina utezanja kraja zida}

\begin{align*}
\varepsilon_{cu2,c} &= 0.0035 + 0.1 \cdot \alpha \cdot \omega_{wd,prov} = 0.0035 + 0.1 \cdot \VAR{d.alfa|n(4)} \cdot \VAR{d.omega_d_prov|n(4)} = \VAR{d.epsilon_cu2c_prov|n(5)} \\
l_{c} &= x_u \cdot \left(1 - \frac{\varepsilon_{cu2}}{\varepsilon_{cu2,c}}\right) = \VAR{(d.xu*100)|n(1)} \cdot \left(1 - \frac{0.0035}{\VAR{d.epsilon_cu2c_prov|n(5)}}\right) = \VAR{(d.lc_req*100)|n(1)} \ \mathrm{cm}
\end{align*}

Kontrola dužine utegnutog elementa:

\begin{align*}
l_c &< h_0 \\
\VAR{(d.lc_req*100)|n(1)} &< \VAR{(m.h_0*100)|n(1)} \quad \VAR{kontrola(d.dovoljno_utegnut)}
\end{align*}

\subsection{Minimalna debljina utegnutog elementa}

Prema EN 1998-1, 5.4.3.4.2(10), debljina $b_w$ utegnutog dela mora biti najmanje 20 cm, a u zavisnosti od dužine utegnutog dela:

\begin{align*}
h_0 &> \max(2 \cdot b_w;\ 0.2 \cdot l_w) = \max(\VAR{(2*m.b*100)|n(1)};\ \VAR{(0.2*m.h*100)|n(1)}) \\
\VAR{(m.h_0*100)|n(1)} &> \VAR{([2*m.b, 0.2*m.h]|max*100)|n(1)} \quad \VAR{uslov(d.b0_min_hs10)}
\end{align*}

\BLOCK{if d.b0_min_hs10}
\begin{align*}
b_{w,min} &= \max\left(20;\ \frac{H_s}{10}\right) = \max\left(20;\ \frac{\VAR{(m.Hs*100)|n(1)}}{10}\right) = \VAR{(d.bw_min*100)|n(1)} \ \mathrm{cm}
\end{align*}
\BLOCK{else}
\begin{align*}
b_{w,min} &= \max\left(20;\ \frac{H_s}{15}\right) = \max\left(20;\ \frac{\VAR{(m.Hs*100)|n(1)}}{15}\right) = \VAR{(d.bw_min*100)|n(1)} \ \mathrm{cm}
\end{align*}
\BLOCK{endif}

Kontrola debljine:

\begin{align*}
b_w &\geq b_{w,min} \\
\VAR{(m.b*100)|n(1)} &\geq \VAR{(d.bw_min*100)|n(1)} \quad \VAR{kontrola(d.debljina_ok)}
\end{align*}

\subsection{Zaključak}

\BLOCK{if not d.dovoljno_omega}
Nije obezbeđen dovoljan procenat utezanja ivičnog elementa ($\omega_{wd,prov} = \VAR{d.omega_d_prov|n(4)} < \omega_{wd,req} = \VAR{d.limit_omega|n(4)}$).

\BLOCK{endif}
\BLOCK{if not d.debljina_ok}
\textbf{Debljina zida $b_w = \VAR{(m.b*100)|n(1)}$ cm je manja od potrebne $b_{w,min} = \VAR{(d.bw_min*100)|n(1)}$ cm.}

\BLOCK{endif}
\BLOCK{if d.dovoljno_utegnut}
Kraj zida \VAR{naziv|tex} je dovoljno utegnut: potrebna dužina utezanja $l_c = \VAR{(d.lc_req*100)|n(1)} \ \mathrm{cm}$ je manja od dužine utegnutog elementa $h_0 = \VAR{(m.h_0*100)|n(1)} \ \mathrm{cm}$.
\BLOCK{else}
\textbf{Potrebno je smanjiti stepen utezanja kraja zida ili produžiti ivični element}: $l_c = \VAR{(d.lc_req*100)|n(1)} \ \mathrm{cm} \geq h_0 = \VAR{(m.h_0*100)|n(1)} \ \mathrm{cm}$.
\BLOCK{endif}
\BLOCK{endif}
\BLOCK{endif}

%% ---- KRAJ PRORACUNA ----

\end{document}
